\noindent{This study addresses the optimization of decentralized, informal paratransit networks in developing regions by formulating a data-driven Transit Route Network Design Problem (\ac{TRNDP}) framework and applying it to the empirical context of Iligan City, Philippines. Traditional route optimization models assume centralized coordination, rigid scheduling, and fixed stop topologies, which fail to capture the flexible boarding-alighting dynamics and uncoordinated nature of paratransit systems like jeepneys. To bridge this gap, this study proposes a parallelized hybrid Memetic Genetic Algorithm (\ac{GA}) integrated with Ant Colony Optimization (\ac{ACO})-inspired Lamarckian learning and an agent-based microscopic passenger simulation. Commuter behavior—including walking thresholds, transfer aversion, and boarding costs—is empirically calibrated using local survey data to construct a high-fidelity travel utility model. Candidate route systems are evaluated in parallel on multi-core processors, using Total User Cost (integrating travel, waiting, and transfer penalties, subject to vehicle capacity limits and fleet equity regularizers) as the primary fitness metric. To overcome the computational bottleneck of evaluating thousands of spatial topologies, an epigenetic inheritance mechanism is implemented where offspring networks warm-start on a blended pheromone memory matrix derived from parent routes. A critical operational gap was resolved by dynamically populating fleet allocations across evolved routes using Mohring's square-root rule during evaluation, thereby providing the Lamarckian local search mutator (comprising spatial attraction, redundancy repulsion, and tortuosity pruning operators) with a valid demand-service gap signal. Post-optimization evaluation across multiple temporal demand regimes (peak and off-peak periods) demonstrates a 15--20\% reduction in average passenger commute times compared to stochastic baselines. Furthermore, the optimized networks display significantly minimized spatial demand-service disparity indices ($D(\mathbf{R})$) and improved equity as measured by fleet distribution Gini coefficients. Structural similarity analysis reveals high cross-run reproducibility (mean Jaccard similarity of 0.73 and 2D Wasserstein distance of 0.009) and high temporal adaptability. This work establishes the feasibility of combining agent-based behavioral modeling and hybrid metaheuristics to redesign complex, informal public transport networks in resource-scarce urban environments.}
